% Generated draft: requires mathematical and source review.
\documentclass[11pt]{article}
\usepackage[margin=25mm]{geometry}
\usepackage{fontspec}
\setmainfont{Latin Modern Roman}
\usepackage{amsmath,amssymb,mathtools,unicode-math}
\setmathfont{Latin Modern Math}
\usepackage{enumitem,needspace,xurl,hyperref}
\setlength{\parindent}{0pt}
\setlength{\parskip}{5pt}
\setlength{\emergencystretch}{3em}
\newcommand{\N}{\mathbb{N}}
\newcommand{\Z}{\mathbb{Z}}
\newcommand{\Q}{\mathbb{Q}}
\newcommand{\R}{\mathbb{R}}
\newcommand{\C}{\mathbb{C}}
\newcommand{\F}{\mathbb{F}}
\newcommand{\A}{\mathbb{A}}
\newcommand{\PP}{\mathbb{P}}
\newcommand{\E}{\mathbb{E}}
\newcommand{\eps}{\varepsilon}
\newcommand{\dd}{\,\mathrm{d}}
\begin{document}
\section*{UnsolvedMath ID: 30002741}
\textbf{Source problem code:} catalogue entry
\subsection{Definitions and mathematical statement}
Let \(\mathcal T\) be a finite tile set intended to tile the Euclidean plane \(\R^2\). Definitions of admissibility are relative to its tile shapes, permitted placements and matching rules; these data are not specified in the supplied record. A global admissible tiling covers \(\R^2\) by permitted tiles with pairwise disjoint interiors and satisfies the matching rules. Write \(X_{\mathcal T}\) for the set of such tilings. The source calls \(\mathcal T\) aperiodic. Under the usual tileable interpretation, this requires \(X_{\mathcal T}\ne\varnothing\) and excludes periodic global tilings; the record does not specify the precise convention for periodicity.

A finite locally admissible patch \(P\) is a finite collection of permitted placed tiles with pairwise disjoint interiors satisfying the applicable local matching rules. No connectedness or boundary-shape restriction is imposed here, since the source supplies none; whether such restrictions are intended remains unresolved. The patch is globally extendible if there exists \(x\in X_{\mathcal T}\) containing every placed tile of \(P\). Thus the globally nonextendible patches form
\[
\mathcal B_{\mathcal T}=\{P: P\text{ is a finite locally admissible patch and }\nexists x\in X_{\mathcal T}\text{ with }P\subseteq x\}.
\]
Here containment means containment as a collection of placed tiles.

A sequential assembly history consists of locally admissible patches \(P_0,P_1,\ldots\), for as long as the assembly continues, with
\[
P_{n+1}=P_n\cup\{t_{n+1}\},
\]
where \(t_{n+1}\) is one new permitted placed tile. The seed, permitted attachment locations, selection rule and stopping rule are unspecified. For a history that has been specified, define its first occurrence of a global obstruction by
\[
\tau=\inf\{n:P_n\in\mathcal B_{\mathcal T}\},\qquad \inf\varnothing=\infty.
\]
Once a patch is globally nonextendible, every subsequent patch containing it is globally nonextendible. This obstruction need not coincide with an immediate inability to attach another tile. Also, an infinite sequence of additions need not cover the whole plane without additional growth requirements.

\textbf{Source problem, with unresolved qualifications.} For a finite aperiodic tile set, determine how often globally nonextendible finite patches cause failure during assembly one tile at a time. A precise quantitative version requires a specified assembly model and a definition of failure. If a probability law on assembly histories is supplied, possible quantities include \(\Pr(\tau\le n)\) and \(\Pr(\tau<\infty)\); if a counting or frequency interpretation is intended, its sampling class and normalization must instead be supplied. These are conditional formulations, not choices made by the source. The supplied record does not determine a unique failure rate, a probability space, or whether the intended result concerns every finite aperiodic tile set or a particular family.
\subsection{Short English statement}
The problem asks how often assembling an aperiodic plane tiling one tile at a time produces a finite patch that cannot belong to any complete tiling. The assembly rules, meaning of failure, and probability or frequency measure are unspecified, so a numerical or asymptotic target cannot yet be stated uniquely. This is a statement-only draft requiring human mathematical review.
\subsection{Sources}
\begin{itemize}
\item Local UnsolvedMath registry, supplied complete record: permanent ID \(30002741\), exact source problem code \texttt{catalogue entry}, title \emph{Failure Rates in Sequential Aperiodic Tiling}. The record belongs to the Oberwolfach Reports Open Problems collection. No individual proposer is supplied; the recorded proposed year is 2014.
\item Supplied source citation: \emph{Mini-Workshop: Dynamical versus Diffraction Spectra in the Theory of Quasicrystals} (2015), Oberwolfach Reports, pp.\ 2987--3013, DOI \texttt{10.4171/owr/2014/53}. Original OWR report: \url{https://doi.org/10.4171/owr/2014/53}. The supplied background describes Oberwolfach Reports as documenting mathematical workshops and their open-problem discussions.
\item Supplied literature source: Koch and Radin, aperiodic tiling model, \url{https://doi.org/10.1007/s10955-009-9896-9}. No fuller bibliographic description or specific supporting passage is supplied.
\item Literature assessment dated exactly \texttt{2026-08-21}: the supplied curation labels the problem open and states that aperiodic tile sets are known to contain finite patches that cannot extend to a global tiling, while a general quantitative theory of the frequency or probability of failure in sequential assembly remains open. This is supplied dated curation derived from the OWR record, not an independently verified literature assessment. Its status and cited claims require independent verification.
\item The registry labels the statement \texttt{corrected\_verified} and reports a manual semantic, typographic and LaTeX review of its cleanup. These are supplied editorial labels; they do not establish that the missing assembly model or quantitative target has been specified. No external source was visited for this draft.
\end{itemize}
\end{document}
